\documentclass[11pt]{article}

\usepackage[a4paper,margin=30mm]{geometry}
\usepackage[T1]{fontenc}
\usepackage{lmodern}
\usepackage{microtype}
\usepackage{amsmath,amssymb,amsthm,mathtools}
\usepackage{booktabs}
\usepackage{enumitem}
\usepackage{xcolor}
\usepackage[hidelinks]{hyperref}

\newtheorem{theorem}{Theorem}[section]
\newtheorem{proposition}[theorem]{Proposition}
\newtheorem{lemma}[theorem]{Lemma}
\newtheorem{corollary}[theorem]{Corollary}
\theoremstyle{definition}
\newtheorem{definition}[theorem]{Definition}
\newtheorem{remark}[theorem]{Remark}

\newcommand{\R}{\mathbb R}
\newcommand{\C}{\mathbb C}
\newcommand{\Z}{\mathbb Z}
\newcommand{\N}{\mathbb N}
\newcommand{\qform}{\mathfrak q}
\newcommand{\Ih}{\mathcal I_h}
\newcommand{\Lh}{L_h}
\newcommand{\supp}{\operatorname{supp}}
\newcommand{\Res}{\operatorname*{Res}}

\title{\textbf{Reflected-Packet Interactions in the Semilocal Weil Form}\\
\large An Exact Zero Expansion and Unconditional Infinite Oscillation}
\author{LeonardSEO}
\date{Version 1.0.0 --- 23 July 2026}

\begin{document}
\maketitle

\begin{abstract}
We freeze one explicit real, even, compactly supported smooth packet and
compute the interaction of two widely separated reflected translates in the
complete semilocal Weil preform.  The pole and all-prime-power sectors combine
into a Stieltjes pairing against \(d\psi-dx\).  Its canonical bilateral
Laplace transform is
\[
L_h(s)\left\{
\frac{\zeta'}{\zeta}\left(s+\frac12\right)
+\frac{1}{s-\frac12}
\right\}.
\]
A contour displacement gives an absolutely convergent expansion over all
nontrivial and trivial zeros.  The archimedean term and the remaining pole
term cancel the complete trivial-zero series, leaving a zero-only expansion
for the full reflected cross interaction.  The packet coefficient vanishes
exactly on the lattice \(4\pi i\Z\setminus\{0\}\).  Conrey's unconditional
positive-proportion theorem therefore supplies infinitely many visible
critical-line poles.  A Laplace-transform form of Landau's theorem then
excludes either eventual sign.  Consequently both the arithmetic interaction
and the even-minus-odd energy splitting take both strict signs arbitrarily
far out.
\end{abstract}

\begin{center}
\fbox{\begin{minipage}{0.91\textwidth}
\textbf{Scope.} This is not a proof of the Riemann hypothesis.  It is a
fixed-packet theorem.  It does not identify, approximate, or determine the
parity of an actual semilocal ground state.  Numerical computations in the
repository are diagnostics only and are not premises of the analytic proof.
\end{minipage}}
\end{center}

\tableofcontents

\section{Introduction and main result}

The Weil explicit formula balances four sectors: an archimedean
distribution, two residual pole terms, the full prime-power measure, and a
spectral sum over the nontrivial zeros.  We study a particularly rigid test:
the interaction of two reflected translates of one fixed packet.  The
purpose is arithmetic.  No spectral ground-state or variational assertion is
made.

Let \(h\in C_c^\infty(\R)\) be the packet constructed in
Section~\ref{sec:packet}.  For \(y>1/2\), set
\[
h_y^{\mathrm R}(x)=h(x-y),\qquad
h_y^{\mathrm L}(x)=h(x+y),
\]
and
\[
h_y^\pm=\frac{h_y^{\mathrm R}\pm h_y^{\mathrm L}}{\sqrt2}.
\]
The supports of \(h_y^{\mathrm R}\) and \(h_y^{\mathrm L}\) are disjoint, so
both \(h_y^\pm\) are normalized.  The energy splitting is
\[
\Delta_h(y)=\qform[h_y^+]-\qform[h_y^-].
\]

Write
\[
K_h=h*h,\qquad
\Lh(w)=\int_{-1}^{1}K_h(r)e^{wr}\,dr
\]
and define the all-prime-power error measure
\[
dE(x)=\sum_{n\ge2}\Lambda(n)\delta_n-dx.
\]
The centered arithmetic interaction is
\[
\Ih(y)=
-\int_0^\infty x^{-1/2}K_h(2y-\log x)\,dE(x).
\]

\begin{theorem}[Exact expansion and unconditional oscillation]
\label{thm:main}
For every \(y>1/2\),
\begin{align}
\Ih(y)
={}&
\sum_{\rho}m_\rho
\Lh\left(\rho-\frac12\right)e^{2(\rho-1/2)y}
\notag\\
&+
\sum_{n=1}^{\infty}
\Lh\left(2n+\frac12\right)e^{-(4n+1)y},
\label{eq:I-expansion}
\end{align}
where \(\rho\) ranges over the distinct nontrivial zeros of \(\zeta\) and
\(m_\rho\) is the multiplicity.  Both series converge absolutely, and the
nontrivial-zero series may be differentiated term by term any fixed number
of times locally uniformly in \(y\).

The archimedean and residual pole terms cancel the second series exactly.
Consequently
\begin{equation}
\qform(h_y^{\mathrm R},h_y^{\mathrm L})
=
\sum_\rho m_\rho
\Lh\left(\rho-\frac12\right)e^{2(\rho-1/2)y},
\label{eq:zero-only}
\end{equation}
and
\begin{equation}
\Delta_h(y)=
2\sum_\rho m_\rho
\Lh\left(\rho-\frac12\right)e^{2(\rho-1/2)y}.
\label{eq:delta-zero-only}
\end{equation}
The coefficient at \(\rho\) is
\[
A_h(\rho)=m_\rho\Lh\left(\rho-\frac12\right),
\]
and it vanishes exactly if
\[
\rho=\frac12+4\pi i k,\qquad k\in\Z\setminus\{0\}.
\]

Unconditionally, for every \(Y>0\), there exist
\(y_+,y_-,u_+,u_->Y\) such that
\[
\Ih(y_+)>0,\qquad \Ih(y_-)<0,
\]
and
\[
\Delta_h(u_+)>0,\qquad \Delta_h(u_-)<0.
\]
Thus each function has infinitely many sign-crossing intervals.
\end{theorem}

The sign convention has a direct but limited interpretation:
\(\Delta_h(y)<0\) means that the even packet has lower energy than the odd
packet at that separation, while \(\Delta_h(y)>0\) means the reverse.
Theorem~\ref{thm:main} therefore proves that neither parity is eventually
preferred for this one packet.

The proof has four independent components.  First, exact polarization fixes
the factor \(2\).  Second, the pole and prime-power main terms are combined
before estimation.  Third, a contour displacement produces
\eqref{eq:I-expansion}, including every trivial zero.  Finally, an exact
packet-transform zero set and Conrey's theorem produce visible nonreal poles;
Landau's theorem then rules out either eventual sign.

\section{Frozen packet and semilocal normalization}
\label{sec:packet}

\subsection{Transform and polarization conventions}

The inner product on \(L^2(\R,dx)\) is antilinear in the first variable:
\[
\langle\Phi,\Psi\rangle
=\int_{\R}\overline{\Phi(x)}\Psi(x)\,dx.
\]
Our Fourier transform and inverse normalization are
\[
\widehat\Phi(z)=\int_{\R}\Phi(x)e^{-izx}\,dx,
\qquad
\langle\Phi,\Psi\rangle
=\frac{1}{2\pi}\int_{\R}
\overline{\widehat\Phi(t)}\widehat\Psi(t)\,dt.
\]
Translation \(T_y\Phi(x)=\Phi(x-y)\) obeys
\[
\widehat{T_y\Phi}(t)=e^{-iyt}\widehat\Phi(t).
\]
For multiplicative tests we use
\[
\widetilde f(s)=\int_0^\infty f(u)u^{s-1}\,du.
\]

For \(\Phi,\Psi\in C_c^\infty(\R)\), define
\[
C_{\Phi,\Psi}(r)
=\int_{\R}\overline{\Phi(x-r)}\Psi(x)\,dx.
\]

\subsection{Complete semilocal Weil preform}

The source normalization reconstructed from
\cite{ConnesConsani2023,CCM2025} is
\begin{align}
\qform(\Phi,\Psi)
={}&
\frac{1}{2\pi}\int_{\R}
m_\infty(t)\,
\overline{\widehat\Phi(t)}\widehat\Psi(t)\,dt
\notag\\
&+
\overline{\widehat\Phi(-i/2)}\widehat\Psi(i/2)
+
\overline{\widehat\Phi(i/2)}\widehat\Psi(-i/2)
\notag\\
&-
\sum_{\substack{p\ {\rm prime}\\m\ge1}}
\frac{\log p}{p^{m/2}}
\left\{
C_{\Phi,\Psi}(m\log p)
+
C_{\Phi,\Psi}(-m\log p)
\right\},
\label{eq:preform}
\end{align}
where
\[
m_\infty(t)
=
\Re\frac{\Gamma'}{\Gamma}
\left(\frac14+\frac{it}{2}\right)-\log\pi.
\]
Every prime power \(p^m\), rather than only every prime, appears with exact
coefficient
\[
-(\log p)p^{-m/2}
=-\Lambda(p^m)(p^m)^{-1/2}.
\]
For a finite semilocal interval \([-\log\lambda,\log\lambda]\), the cutoff is
\[
p^m\le\lambda^2.
\]
The endpoint is inclusive and carries no half weight.

The equivalent real-place distribution is
\begin{align}
W_{\R}(C)
={}&
(\log 4\pi+\gamma_E)C(0)
\notag\\
&+
\int_0^\infty
\frac{
e^{r/2}\{C(r)+C(-r)\}-2C(0)
}{
e^r-e^{-r}
}\,dr,
\label{eq:real-place}
\end{align}
and the archimedean contribution to \eqref{eq:preform} is
\(-W_{\R}\).  This fixes the scalar, subtraction, and sign.

\subsection{The dyadic infinite-convolution packet}

For \(a_n=2^{-n-1}\), \(n\ge1\), put
\[
b_{a_n}(x)=\frac{1}{2a_n}\mathbf1_{[-a_n,a_n]}(x)
\]
and define
\[
g=*_{n\ge1}b_{a_n},\qquad
h=\frac{g}{c_g},\qquad c_g=\lVert g\rVert_2.
\]
Since \(\sum_{n\ge1}a_n=1/2\), the infinite convolution is the density
of an absolutely convergent sum of independent compactly supported random
variables.  Moreover
\[
g\in C_c^\infty(\R),\qquad
\supp g=[-1/2,1/2],
\]
and \(g\) is real, even, nonnegative, and strictly positive in the
interior.  Smoothness follows from the super-polynomial Fourier decay
obtained by bounding any fixed number of sinc factors.

The transform is
\[
\widehat h(z)
=c_g^{-1}
\prod_{n\ge1}\frac{\sin(a_nz)}{a_nz}.
\]
The product converges locally uniformly, and
\[
Z(\widehat h)=4\pi\Z\setminus\{0\}.
\]

Because \(h\) is even,
\[
K_h(r)=\int_\R h(x-r)h(x)\,dx=(h*h)(r).
\]
It is real, even, nonnegative, smooth, supported exactly on \([-1,1]\),
positive on \((-1,1)\), and \(K_h(0)=1\).  Its bilateral Laplace transform
is
\begin{equation}
\Lh(w)
=c_g^{-2}
\prod_{n\ge1}
\left(\frac{\sinh(a_nw)}{a_nw}\right)^2.
\label{eq:L-product}
\end{equation}
This is entire and even.  Since
\[
\frac{\sinh(a_nw)}{a_nw}=1+O(a_n^2)
\]
uniformly on compact \(w\)-sets and \(\sum a_n^2<\infty\), the infinite
tail has no extra zeros.  Hence
\begin{equation}
Z(\Lh)=4\pi i\Z\setminus\{0\}.
\label{eq:L-zero-set}
\end{equation}
At \(4\pi i k\ne0\), the multiplicity is
\[
2\bigl(v_2(|k|)+1\bigr).
\]

\section{Polarization and exact prime--pole interaction}
\label{sec:interaction}

For \(y>1/2\), the two translates have disjoint support.  Hermitian
polarization, with the first form variable antilinear, gives
\begin{align*}
\qform[h_y^+]
&=\frac12\{\qform[h_y^{\mathrm R}]+\qform[h_y^{\mathrm L}]\}
+\Re\qform(h_y^{\mathrm R},h_y^{\mathrm L}),\\
\qform[h_y^-]
&=\frac12\{\qform[h_y^{\mathrm R}]+\qform[h_y^{\mathrm L}]\}
-\Re\qform(h_y^{\mathrm R},h_y^{\mathrm L}).
\end{align*}
Inversion symmetry makes the cross term real.  Thus
\begin{equation}
\boxed{\Delta_h(y)=2\qform(h_y^{\mathrm R},h_y^{\mathrm L}).}
\label{eq:polarization}
\end{equation}

Let
\[
M_h=\widehat h(i/2)=\widehat h(-i/2)
=\int_\R h(x)e^{x/2}\,dx>0.
\]
The two pole terms in \eqref{eq:preform} contribute exactly
\begin{equation}
M_h^2(e^y+e^{-y}).
\label{eq:pole-cross}
\end{equation}
The prime-power cross term is
\begin{equation}
-
\sum_{p,m\ge1}
\frac{\log p}{p^{m/2}}
K_h(2y-m\log p).
\label{eq:prime-cross}
\end{equation}
Because \(K_h\) is supported on \([-1,1]\), only
\[
2y-1<m\log p<2y+1
\]
contributes.  Equality is zero because \(K_h(\pm1)=0\).  This is a support
fact, not an endpoint half-weight convention.

Define the full Chebyshev measure
\[
d\Psi=\sum_{n\ge2}\Lambda(n)\delta_n,\qquad
\psi(X)=\sum_{n\le X}\Lambda(n),
\]
and
\[
dE=d\Psi-dx.
\]
This is the prime-power error distribution associated with
\(\psi(X)-X\), not the primes-only distribution associated with
\(\theta(X)-X\).

The continuous term can be evaluated without approximation:
\begin{align}
\int_0^\infty
x^{-1/2}K_h(2y-\log x)\,dx
&=
e^y\int_{-1}^{1}K_h(r)e^{-r/2}\,dr
\notag\\
&=e^yM_h^2.
\label{eq:continuous-main}
\end{align}
Therefore
\begin{align}
\Ih(y)
&=
e^yM_h^2
-
\sum_{n\ge2}
\frac{\Lambda(n)}{\sqrt n}
K_h(2y-\log n),
\label{eq:I-physical}
\end{align}
and the exact cross interaction is
\begin{equation}
\qform(h_y^{\mathrm R},h_y^{\mathrm L})
=
A_h^\infty(y)+e^{-y}M_h^2+\Ih(y),
\label{eq:Q-decomposition}
\end{equation}
where
\[
A_h^\infty(y)
=
\frac{1}{2\pi}
\int_\R m_\infty(t)e^{2iyt}|\widehat h(t)|^2\,dt.
\]
Equation~\eqref{eq:Q-decomposition} is the exact prime--pole cancellation:
the \(e^yM_h^2\) pole main term has been combined with the continuous part
of \(dE\) before any estimate is made.

For a packet realized inside
\([-\log\lambda,\log\lambda]\), the geometric condition is
\[
y+\frac12\le\log\lambda.
\]
Every contributing prime power then automatically satisfies the inclusive
semilocal cutoff \(p^m\le\lambda^2\).

\section{Canonical transform and exact zero expansion}
\label{sec:explicit}

\subsection{The base-1 extension}

The uncut Mellin transform of \(dE\) is not defined in a right half-plane,
because \(dE=-dx\) on \((0,1)\).  Introduce only for the transform
\[
J_h(t)=
-\int_{[1,\infty)}
x^{-1/2}K_h(t-\log x)\,dE(x).
\]
If
\[
E_0(x)=\psi(x)-x,\qquad
E_1(x)=\psi(x)-(x-1)\quad(x\ge1),
\]
then \(E_0(1)=-1\), \(E_1(1)=0\), and distributionally
\[
d(\mathbf1_{[1,\infty)}E_0)
=\mathbf1_{[1,\infty)}dE-\delta_1,
\]
whereas
\[
d(\mathbf1_{[1,\infty)}E_1)
=\mathbf1_{[1,\infty)}dE.
\]
Thus \(J_h\) uses the endpoint-normalized primitive \(E_1\) and contains no
spurious \(K_h(t)\) boundary term.  It satisfies
\[
J_h(t)=0\quad(t\le-1),\qquad
J_h(t)=\Ih(t/2)\quad(t>1).
\]

\subsection{Bilateral Laplace transform}

For \(\Re s>1/2\), absolute Euler-product convergence and Tonelli's theorem
give
\begin{align}
\mathcal B_h(s)
&=
\int_\R J_h(t)e^{-st}\,dt
\notag\\
&=
\Lh(s)
\left\{
\frac{\zeta'}{\zeta}\left(s+\frac12\right)
+\frac{1}{s-\frac12}
\right\}.
\label{eq:B-transform}
\end{align}
At \(s=1/2\),
\[
\frac{\zeta'}{\zeta}\left(s+\frac12\right)
=-\frac{1}{s-1/2}+\gamma_E+O(s-1/2),
\]
so the apparent pole is removable and
\[
\mathcal B_h(1/2)=\gamma_E\Lh(1/2)=\gamma_E M_h^2.
\]
The remaining poles are
\[
s=\rho-\frac12
\]
for nontrivial zeros, with residue
\[
m_\rho\Lh\left(\rho-\frac12\right),
\]
and
\[
s=-2n-\frac12,\qquad n\ge1,
\]
for trivial zeros, with residue
\[
\Lh\left(-2n-\frac12\right).
\]
A multiple zero changes the residue by its multiplicity; it does not create
a derivative term.

\subsection{Contour displacement}

For \(c>1/2\), Fourier inversion on the line \(\Re s=c\) gives
\[
J_h(t)
=
\frac{1}{2\pi i}\int_{c-i\infty}^{c+i\infty}
e^{st}\Lh(s)
\left\{
\frac{\zeta'}{\zeta}\left(s+\frac12\right)
+\frac{1}{s-\frac12}
\right\}ds.
\]
The integral is absolutely convergent because, uniformly in every fixed
vertical strip,
\[
\Lh(\sigma+i\tau)=O_N((1+|\tau|)^{-N})
\]
for every \(N\).  Choose heights avoiding zero ordinates by
\(\gg1/\log T\); on those horizontal segments
\(\zeta'/\zeta=O(\log^2T)\), so the horizontal integrals vanish.

Move the left side to \(\Re s=-2N-1\).  Repeated integration by parts and
\(\supp K_h=[-1,1]\) give, for \(M\ge3\),
\[
|\Lh(-2N-1+i\tau)|
\le
C_Me^{2N+1}(1+N)^M(1+|\tau|)^{-M}.
\]
The remaining vertical integral is consequently
\[
O_{h,M}\left(
e^{-(2N+1)(t-1)}(1+N)^M\log(N+2)
\right),
\]
which tends to zero when \(t>1\).  Summing the crossed residues yields
\begin{align}
J_h(t)
={}&
\sum_\rho m_\rho
\Lh\left(\rho-\frac12\right)
e^{(\rho-1/2)t}
\notag\\
&+
\sum_{n=1}^{\infty}
\Lh\left(-2n-\frac12\right)e^{(-2n-1/2)t}.
\label{eq:J-explicit}
\end{align}
Since \(\Lh\) is even, \(t=2y\) proves
\eqref{eq:I-expansion}.

The trivial-zero sum is absolutely convergent for \(t>1\).  The zero-counting
bound \(N_\zeta(T)=O(T\log T)\), combined with rapid vertical decay of
\(\Lh\), proves absolute and locally uniform convergence of the nontrivial
zero sum after any fixed number of real derivatives.

\subsection{Archimedean cancellation}

For the reflected cross correlation, \(C(0)=0\), \(C(r)=0\) for \(r\ge0\),
and \(C(-r)=K_h(2y-r)\).  Equation~\eqref{eq:real-place} gives
\[
A_h^\infty(y)
=-\int_0^\infty
\frac{e^{r/2}}{e^r-e^{-r}}K_h(2y-r)\,dr.
\]
Using
\[
\frac{e^{r/2}}{e^r-e^{-r}}
=\sum_{n=0}^{\infty}e^{-(2n+1/2)r},
\]
Tonelli's theorem and \(u=2y-r\) give
\[
A_h^\infty(y)
=-\sum_{n=0}^{\infty}
\Lh\left(2n+\frac12\right)e^{-(4n+1)y}.
\]
The \(n=0\) term is \(-M_h^2e^{-y}\); the terms \(n\ge1\) cancel the
complete trivial-zero series in \eqref{eq:I-expansion}.  Combining with
\eqref{eq:Q-decomposition} proves the zero-only formulas
\eqref{eq:zero-only} and \eqref{eq:delta-zero-only}.

\section{Visible zero coefficients and infinite oscillation}
\label{sec:oscillation}

\subsection{Coefficient visibility}

From \eqref{eq:L-product}, the \(n\)-th factor vanishes precisely at
\[
w=\pi i\ell\,2^{n+1},\qquad \ell\in\Z\setminus\{0\}.
\]
The locally uniform nonvanishing of the tail away from factor zeros proves
\eqref{eq:L-zero-set}.  Thus every off-critical-line zero is visible.  A
critical-line zero is invisible only if its ordinate is in
\[
4\pi\Z\setminus\{0\}.
\]

Conrey's Theorem~1 in \cite{Conrey1989} proves that a positive proportion
of all nontrivial zeros are simple and lie on the critical line.  Hence
there are \(\gg T\log T\) distinct simple critical-line zeros below height
\(T\).  The invisible lattice has only \(O(T)\) ordinates in that range,
and a simple zero occupies at most one such ordinate.  It follows
unconditionally that infinitely many critical-line zeros are visible for
the frozen packet.

\subsection{A Laplace-transform Landau lemma}

\begin{lemma}[Landau]
\label{lem:landau}
Let \(T_0\ge0\), and let
\(f:[T_0,\infty)\to[0,\infty)\) be locally integrable.  Suppose its Laplace
integral converges on a nonempty right half-line, and define
\[
F(s)=\int_{T_0}^{\infty}f(t)e^{-st}\,dt,
\]
\[
\sigma_c=
\inf\left\{
\sigma\in\R:
\int_{T_0}^{\infty}f(t)e^{-\sigma t}\,dt<\infty
\right\}.
\]
If \(\sigma_c\) is finite, then \(F\) is singular at the real point
\(s=\sigma_c\).
\end{lemma}

\begin{proof}
Assume that \(F\) is holomorphic near \(\sigma_c\).  Choose
\(\sigma_1>\sigma_c\) and \(r>0\) so that
\(\sigma_1-r<\sigma_c\) while the Taylor disk from \(\sigma_1\) to
\(\sigma_1-r\) remains inside the union of this neighborhood and
\(\Re s>\sigma_c\).  For a small \(\varepsilon>0\) with
\(\sigma_1-\varepsilon>\sigma_c\), the inequality
\(t^n\le C_{n,\varepsilon}e^{\varepsilon t}\) justifies differentiation:
\[
(-1)^nF^{(n)}(\sigma_1)
=\int_{T_0}^{\infty}t^nf(t)e^{-\sigma_1t}\,dt\ge0.
\]
Taylor's theorem and monotone convergence give
\begin{align*}
F(\sigma_1-r)
&=
\sum_{n=0}^{\infty}
\frac{(-r)^nF^{(n)}(\sigma_1)}{n!}\\
&=
\int_{T_0}^{\infty}
f(t)e^{-\sigma_1t}
\sum_{n=0}^{\infty}\frac{(rt)^n}{n!}\,dt\\
&=
\int_{T_0}^{\infty}
f(t)e^{-(\sigma_1-r)t}\,dt<\infty,
\end{align*}
contradicting the definition of \(\sigma_c\).
\end{proof}

\subsection{Oscillation of the arithmetic interaction}

The direct prime-power window gives the elementary bound
\[
|\Ih(t/2)|=O_h((1+t)e^{t/2}),\qquad t>1.
\]
Indeed, one may use
\(\sum_{n\le X}\Lambda(n)\le X\log X\).  Thus the tail transform has a
finite abscissa of convergence no greater than \(1/2\).

At each visible critical-line zero \(\rho=1/2+i\gamma\), the meromorphic
continuation \eqref{eq:B-transform} has a genuine pole at \(s=i\gamma\).
Suppose \(\Ih(t/2)\) were eventually nonnegative.  Discarding the finite
initial interval changes its transform by an entire function.  The visible
pole rules out an abscissa below \(0\), while absolute convergence gives
\[
0\le\sigma_c\le\frac12.
\]
But \eqref{eq:B-transform} is holomorphic at every real point of this
interval: the pole at \(1/2\) cancels, and
\(\zeta(u)\ne0\) for real \(1/2\le u<1\).  Lemma~\ref{lem:landau} gives a
contradiction.  Applying the same argument to \(-\Ih(t/2)\) excludes
eventual nonpositivity.

\subsection{Oscillation of the complete splitting}

Put \(\mathcal Q_h(t)=\qform(h_{t/2}^{\mathrm R},
h_{t/2}^{\mathrm L})\).  The zero-only formula gives
\[
\mathcal Q_h(t)=
\sum_\rho a_\rho e^{(\rho-1/2)t},
\qquad
a_\rho=m_\rho\Lh\left(\rho-\frac12\right).
\]
For every \(T_0>1\),
\[
\int_{T_0}^{\infty}\mathcal Q_h(t)e^{-st}\,dt
=
\sum_\rho
\frac{
a_\rho e^{(\rho-1/2-s)T_0}
}{
s-(\rho-1/2)
}.
\]
Rapid coefficient decay makes the right side normally convergent away
from its displayed poles.  It has genuine nonreal poles at every visible
critical zero and no real poles.  The same Landau argument, applied
separately to \(\mathcal Q_h\) and to \(-\mathcal Q_h\), excludes either
eventual sign.  Equation~\eqref{eq:polarization} transfers the conclusion to
\(\Delta_h\).

Both \(\Ih\) and \(\Delta_h\) are continuous.  Alternating positive and
negative points can be chosen beyond increasing bounds, producing
infinitely many disjoint sign-crossing intervals.  This does not assert
that the intervening zeros are simple or transverse.

\section{Numerical diagnostics and finite certificates}
\label{sec:numerics}

No numerical statement in this section is used in
Theorem~\ref{thm:main}.  The analytic nonvanishing input comes from the exact
zero set \eqref{eq:L-zero-set} and Conrey's theorem, not from computed zeta
zeros.

\subsection{Implementation A: direct prime-power side}

The first implementation evaluates finite dyadic convolutions by an exact
nonuniform B-spline identity.  It enumerates every \(p^m\) in
\[
|2y-\log(p^m)|<1
\]
and records the triple \((p,m,p^m)\) and the exact coefficient
\((\log p)/\sqrt{p^m}\).  Packet truncations \(N=12,14,16\) are retained
separately.  An exact-rational computation certifies
\[
0.063637725597
<
\Ih(1/2)
<
0.088006928482.
\]
The finite window at \(y=1/2\) contains exactly
\[
2,\ 3,\ 4,\ 5,\ 7.
\]
The certificate includes the omitted infinite-convolution tail and does not
replace the infinite oscillation proof.

\subsection{Implementation B: independent Fourier and zero sides}

The second implementation computes \(K_h\) by Fourier inversion of the sinc
product on its direct side.  Independently, it evaluates the truncated
explicit formula through the product \eqref{eq:L-product} and zeta zeros.
At \(y=1.25,1.75,2.25\), the direct calculations retain respectively
\(15,26,50\) prime powers.  The two implementations agree to between
\(5.2\times10^{-13}\) and \(1.4\times10^{-11}\), inside their declared
diagnostic budgets.

At 220-bit Arb precision, selected coefficients at zero indices \(1,2,4\)
are certified to have strictly positive real intervals after enclosing the
infinite product tail.  The omitted nontrivial-zero sum is not certified;
therefore this computation makes no eventual-sign claim.

\begin{table}[ht]
\centering
\begin{tabular}{@{}rrr@{}}
\toprule
\(y\) & retained prime powers & \(|\Ih^{(A)}-\Ih^{(B)}|\)\\
\midrule
1.25 & 15 & \(5.13\times10^{-13}\)\\
1.75 & 26 & \(8.25\times10^{-12}\)\\
2.25 & 50 & \(1.32\times10^{-11}\)\\
\bottomrule
\end{tabular}
\caption{Cross-implementation diagnostics.}
\end{table}

Machine-readable JSON and CSV files in the repository contain every retained
prime power, including composite prime powers, rather than a primes-only
projection.

\section{Logical scope and independent audit}

The analytic chain was reconstructed independently multiple times and
subjected to separate circularity audits.  The audited dependencies are:
\begin{itemize}[leftmargin=2em]
\item the frozen semilocal normalization
      \eqref{eq:preform};
\item elementary properties of the fixed infinite convolution;
\item absolute convergence of the Euler product in its initial half-plane;
\item standard meromorphic continuation and the functional equation of
      \(\zeta\);
\item the zero-counting estimate \(N_\zeta(T)=O(T\log T)\);
\item Conrey's unconditional positive-proportion theorem
      \cite{Conrey1989}; and
\item the Landau lemma proved in this paper.
\end{itemize}

The proof does not assume:
\begin{itemize}[leftmargin=2em]
\item the Riemann hypothesis;
\item global simplicity of zeta zeros;
\item linear independence of zero ordinates;
\item Weil positivity;
\item an eventual sign;
\item a numerical zero table;
\item convergence of a determinant;
\item completeness of any spectral realization; or
\item information about semilocal ground states.
\end{itemize}

The exact coefficient transform was fixed before any zero theorem was
applied.  Its invisible lattice was derived from the packet alone, so the
visible poles are not inserted by normalization.  The two numerical
implementations can be deleted without changing any step of the proof.

\begin{corollary}[Fixed-packet no eventual parity]
For the frozen packet,
\[
\Delta_h(y)\le0\quad\text{for all sufficiently large }y
\]
and
\[
\Delta_h(y)\ge0\quad\text{for all sufficiently large }y
\]
are both false.
\end{corollary}

This corollary is deliberately not a theorem about the parity of a true
ground state.  Such a statement would require uniform estimates over a
packet class and a variational comparison controlling the magnitude of the
interaction.  Neither is claimed here.


\begin{thebibliography}{9}

\bibitem{ConnesConsani2023}
A.~Connes and C.~Consani,
\emph{Spectral triples and \(\zeta\)-cycles},
Enseign. Math. \textbf{69} (2023), 93--148.
\href{https://doi.org/10.4171/LEM/1049}{doi:10.4171/LEM/1049}.

\bibitem{CCM2025}
A.~Connes, C.~Consani, and H.~Moscovici,
\emph{Zeta Spectral Triples},
arXiv:2511.22755 (2025).
\href{https://arxiv.org/abs/2511.22755}{arXiv:2511.22755}.

\bibitem{Conrey1989}
J.~B. Conrey,
\emph{More than two fifths of the zeros of the Riemann zeta function
are on the critical line},
J. Reine Angew. Math. \textbf{399} (1989), 1--26.
\href{https://doi.org/10.1515/crll.1989.399.1}
{doi:10.1515/crll.1989.399.1}.

\end{thebibliography}

\end{document}
